
% !TEX program = xelatex
\documentclass[11pt,a4paper,oneside]{book}
\usepackage{fontspec}
\usepackage[ngerman]{babel}
\usepackage{microtype}
\usepackage{amsmath,amssymb,amsthm,mathtools}
\usepackage{geometry}
\usepackage{enumitem}
\usepackage{booktabs}
\usepackage{array}
\usepackage{xcolor}
\usepackage{hyperref}
\usepackage{titlesec}
\usepackage{fancyhdr}
\usepackage{setspace}
\usepackage{csquotes}
\usepackage{tocloft}
\geometry{left=27mm,right=27mm,top=24mm,bottom=29mm,headheight=15pt}
\setmainfont{Liberation Serif}
\setsansfont{Liberation Sans}
\setmonofont{Liberation Mono}
\hypersetup{colorlinks=true,linkcolor=black,urlcolor=black,citecolor=black,pdfauthor={Ingolf Lohmann},pdftitle={Quantengravitation aus Quantenprinzip und Raumzeitdynamik}}
\setlength{\parindent}{0pt}
\setlength{\parskip}{6pt}
\setstretch{1.08}
\emergencystretch=3em
\sloppy
\pagestyle{fancy}
\fancyhf{}
\fancyhead[L]{Quantengravitation aus Quantenprinzip und Raumzeitdynamik}
\fancyhead[R]{\thepage}
\renewcommand{\headrulewidth}{0.4pt}
\titleformat{\chapter}{\normalfont\LARGE\bfseries}{\thechapter.}{0.75em}{}
\titleformat{\section}{\normalfont\Large\bfseries}{\thesection}{0.6em}{}
\titleformat{\subsection}{\normalfont\large\bfseries}{\thesubsection}{0.6em}{}
\theoremstyle{definition}
\newtheorem{definition}{Definition}[chapter]
\newtheorem{beispiel}[definition]{Beispiel}
\newtheorem{annahme}[definition]{Annahme}
\theoremstyle{plain}
\newtheorem{satz}[definition]{Satz}
\newtheorem{lemma}[definition]{Lemma}
\newtheorem{korollar}[definition]{Korollar}
\theoremstyle{remark}
\newtheorem*{bemerkung}{Bemerkung}
\newenvironment{merkbox}{\begin{quote}\itshape}{\end{quote}}
\newcommand{\R}{\mathbb{R}}
\newcommand{\C}{\mathbb{C}}
\newcommand{\M}{\mathcal{M}}
\newcommand{\Hh}{\mathcal{H}}
\newcommand{\Lag}{\mathcal{L}}
\newcommand{\G}{\mathcal{G}}
\newcommand{\E}{\mathcal{E}}
\newcommand{\dd}{\mathrm{d}}
\newcommand{\order}{\mathcal{O}}
\newcommand{\ket}[1]{\left|#1\right\rangle}
\newcommand{\bra}[1]{\left\langle#1\right|}
\newcommand{\braket}[2]{\left\langle#1\middle|#2\right\rangle}
\begin{document}
\frontmatter
\begin{titlepage}
\centering
\vspace*{18mm}
{\Huge\bfseries Quantengravitation aus Quantenprinzip und Raumzeitdynamik\par}
\vspace{5mm}
{\Large Ein mathematisches Beweisskript für das zweite Semester\par}
\vspace{13mm}
{\large Von der klassischen Raumzeit zur notwendig quantisierten Gravitation\par}
\vfill
{\large q.e.d.\\Ingolf Lohmann\par}
\vspace{12mm}
{\large Fassung V5 -- korrigierte mathematisch-physikalische Fassung\par}
\vspace{6mm}
{\large 2026\par}
\end{titlepage}

\chapter*{Lizenz- und Nutzungsvereinbarung}
\addcontentsline{toc}{chapter}{Lizenz- und Nutzungsvereinbarung}
Autor und Urheber: Ingolf Lohmann.

Dieses Skript ist als Lehr-, Prüf- und Forschungsunterlage bestimmt. Soweit nicht ausdrücklich schriftlich anders vereinbart, gilt für den Text die Lizenz \textbf{CC BY-NC-ND 4.0}. Das bedeutet: Namensnennung ist erforderlich; eine kommerzielle Nutzung ist nicht erlaubt; Bearbeitungen oder abgeleitete Fassungen dürfen nicht ohne gesonderte Zustimmung verbreitet werden.

Die mathematischen und physikalischen Standardbegriffe, Gleichungen und Lehrsätze, die hier verwendet werden, bleiben als allgemein bekannte wissenschaftliche Grundlagen frei. Geschützt ist die konkrete Darstellung, Auswahl, Reihenfolge, Formulierung und didaktische Beweisführung dieses Skripts.

\chapter*{Vorbemerkung}
\addcontentsline{toc}{chapter}{Vorbemerkung}
Dieses Skript verfolgt ein klares Ziel: Es soll für Studierende im zweiten Semester nachvollziehbar machen, weshalb eine konsistente Verbindung von Quantenprinzip und dynamischer Raumzeit nicht bei einer rein klassischen Gravitation stehen bleiben kann.

Der Beweisanspruch dieses Skripts ist präzise: Aus bekannten mathematischen und physikalischen Voraussetzungen wird gezeigt, dass die gravitativen Freiheitsgrade selbst quantisiert werden müssen, sobald sie konsistent mit quantisierter Materie, Superposition, lokaler Energie-Impuls-Erhaltung und dynamischer Raumzeit wechselwirken sollen. Damit wird die Notwendigkeit einer Quantengravitation bewiesen. Nicht bewiesen wird in diesem Skript die endgültige Auswahl einer bestimmten vollständigen Theorie bei Planck-Energien. Das wäre ein stärkerer empirischer und theoretischer Anspruch.

Die Formulierung ist bewusst elementar aufgebaut. Zunächst werden Raumzeit, Felder, Wirkung, Variationsprinzip, Quantenprinzip und Energie-Impuls-Tensor eingeführt. Danach folgt die Beweiskette: klassische Gravitation ist Dynamik der Metrik; Quantenmaterie besitzt Superpositionen des Energie-Impuls-Tensors; eine rein klassische Metrik kann diese Superposition nicht ohne Strukturverlust, Nichtlinearitätsprobleme oder Messwiderspruch tragen; die schwache Feldnäherung zeigt die gravitative Störung als masseloses Spin-2-Feld; Quantisierung dieses Feldes ist die niederenergetisch zwingende Form der Quantengravitation.

\tableofcontents
\mainmatter

\chapter{Das Beweisziel}
\section{Die Leitfrage}
Die Leitfrage lautet nicht: Welche gegenwärtig diskutierte Theorie der Quantengravitation ist endgültig richtig?

Die Leitfrage lautet:

\begin{merkbox}
Warum folgt aus bekannter Mathematik und bekannter Physik, dass Gravitation nicht dauerhaft rein klassisch bleiben kann, sobald Materie quantisiert ist?
\end{merkbox}

Diese Frage ist stärker als eine bloße Vermutung und schwächer als die Behauptung, bereits eine vollständige Theorie aller Planck-Skalen-Phänomene formuliert zu haben.

Sie führt zu einem präzisen Satz:

\begin{satz}[Notwendigkeit der Quantengravitation]
Wenn Materie quantenmechanisch beschrieben wird, wenn Energie und Impuls die Quelle der Gravitation sind, wenn Gravitation dynamische Raumzeitgeometrie ist und wenn Wechselwirkungen lokal, konsistent und ohne widersprüchliche Messstruktur beschrieben werden sollen, dann müssen auch die gravitativen Freiheitsgrade quantisierte Freiheitsgrade besitzen. Eine rein klassische Gravitation ist dann keine vollständige konsistente Endbeschreibung.
\end{satz}

Der Satz sagt nicht, ob die richtige Volltheorie Schleifen, Strings, kausale Mengen, Pfadintegrale, Spin-Netzwerke, Holographie oder eine andere Struktur verwendet. Er sagt: Irgendeine Form quantisierter Gravitation ist notwendig.

\section{Was heißt hier Beweis?}
Ein Beweis in diesem Skript besteht aus vier Teilen:

\begin{enumerate}[label=\arabic*.]
\item Die Prämissen werden offen genannt.
\item Jeder verwendete mathematische Schritt wird definiert.
\item Jede physikalische Schlussfolgerung wird an eine bekannte Gleichung oder ein bekanntes Prinzip gebunden.
\item Der Schluss wird auf seine Reichweite begrenzt.
\end{enumerate}

Ein Beweis ist nicht dadurch stärker, dass er mehr behauptet. Er ist stärker, wenn seine Voraussetzungen, Folgerungen und Grenzen klar sind.

\section{Die fünf Grundpfeiler}
Die Herleitung verwendet fünf allgemein bekannte Grundpfeiler:

\begin{enumerate}[label=\textbf{P\arabic*.}]
\item Raumzeit wird in der Allgemeinen Relativitätstheorie durch eine Lorentz-Mannigfaltigkeit mit Metrik beschrieben.
\item Materie und Strahlung besitzen einen Energie-Impuls-Tensor.
\item Die Metrik ist dynamisch und wird durch Energie und Impuls beeinflusst.
\item Quantenmechanische Zustände können Superpositionen physikalisch unterscheidbarer Alternativen bilden.
\item Wechselwirkungen müssen konsistent mit Erhaltungssätzen, Messstatistik und lokaler Dynamik sein.
\end{enumerate}

Aus diesen Pfeilern wird die Beweiskette aufgebaut.

\chapter{Mathematische Grundlagen}
\section{Mannigfaltigkeiten}
\begin{definition}[Mannigfaltigkeit]
Eine vierdimensionale glatte Mannigfaltigkeit $\M$ ist ein Raum, der lokal wie $\R^4$ aussieht und auf dem differenzierbare Koordinatenwechsel möglich sind.
\end{definition}

Die physikalische Idee lautet: Beobachter können lokal Koordinaten verwenden, aber keine Koordinate ist die Wirklichkeit selbst. Koordinaten sind Beschreibungen. Die Geometrie muss koordinatenunabhängig formuliert werden.

\section{Tensoren}
Ein Tensor ist ein geometrisches Objekt, dessen Komponenten sich bei Koordinatenwechseln nach festen Regeln transformieren. Tensoren erlauben es, physikalische Gesetze koordinatenunabhängig zu formulieren.

Ein Vektor $V$ besitzt Komponenten $V^\mu$. Eine kovariante Größe besitzt Komponenten $\omega_\mu$. Ein Tensor zweiter Stufe besitzt zum Beispiel Komponenten $T_{\mu\nu}$ oder $T^{\mu\nu}$.

\section{Die Metrik}
\begin{definition}[Metrik]
Eine Lorentz-Metrik $g_{\mu\nu}$ auf einer vierdimensionalen Mannigfaltigkeit ist ein symmetrisches, nicht entartetes Tensorfeld mit Signatur $(-,+,+,+)$ oder $(+,-,-,-)$.
\end{definition}

Die Metrik bestimmt das Linienelement
\[
\dd s^2 = g_{\mu\nu}\,\dd x^\mu\dd x^\nu.
\]

Sie bestimmt, welche Abstände zeitartig, lichtartig oder raumartig sind. Damit beschreibt sie nicht bloß Abstände, sondern auch die Kausalstruktur.

\section{Kausalstruktur}
Für einen Tangentialvektor $v^\mu$ gilt bei Signatur $(-,+,+,+)$:
\[
g_{\mu\nu}v^\mu v^\nu < 0 \quad \text{zeitartig},
\]
\[
g_{\mu\nu}v^\mu v^\nu = 0 \quad \text{lichtartig},
\]
\[
g_{\mu\nu}v^\mu v^\nu > 0 \quad \text{raumartig}.
\]

Damit legt die Metrik fest, was kausal verbunden sein kann. Wer die Metrik dynamisch macht, macht auch die Kausalstruktur dynamisch.

\chapter{Klassische Gravitation als Geometrie}
\section{Geodäten}
Freie Teilchen bewegen sich in der Allgemeinen Relativitätstheorie entlang Geodäten. Eine Geodäte ist eine Kurve $x^\mu(\lambda)$, die lokal die natürliche Gerade der gekrümmten Geometrie darstellt.

Die Geodätengleichung lautet:
\[
\frac{\dd^2 x^\mu}{\dd\lambda^2}
+\Gamma^\mu_{\alpha\beta}\frac{\dd x^\alpha}{\dd\lambda}
\frac{\dd x^\beta}{\dd\lambda}=0.
\]

Die Christoffelsymbole sind
\[
\Gamma^\mu_{\alpha\beta}=\frac12 g^{\mu\nu}\left(\partial_\alpha g_{\nu\beta}+\partial_\beta g_{\nu\alpha}-\partial_\nu g_{\alpha\beta}\right).
\]

Sie hängen von der Metrik ab. Also bestimmt die Metrik die Bewegung freier Körper.

\section{Krümmung}
Die Krümmung wird durch den Riemannschen Krümmungstensor beschrieben:
\[
R^\rho{}_{\sigma\mu\nu}
=\partial_\mu\Gamma^\rho_{\nu\sigma}
-\partial_\nu\Gamma^\rho_{\mu\sigma}
+\Gamma^\rho_{\mu\lambda}\Gamma^\lambda_{\nu\sigma}
-\Gamma^\rho_{\nu\lambda}\Gamma^\lambda_{\mu\sigma}.
\]

Aus ihm erhält man den Ricci-Tensor
\[
R_{\mu\nu}=R^\rho{}_{\mu\rho\nu},
\]
und den Ricci-Skalar
\[
R=g^{\mu\nu}R_{\mu\nu}.
\]

\section{Einstein-Tensor}
Der Einstein-Tensor lautet
\[
G_{\mu\nu}=R_{\mu\nu}-\frac12 Rg_{\mu\nu}.
\]

Er erfüllt die geometrische Identität
\[
\nabla^\mu G_{\mu\nu}=0.
\]

Diese Identität ist entscheidend, weil sie zur lokalen Energie-Impuls-Erhaltung passt.

\section{Einsteins Feldgleichung}
Die klassische Feldgleichung lautet
\[
G_{\mu\nu}+\Lambda g_{\mu\nu}=\frac{8\pi G}{c^4}T_{\mu\nu}.
\]

Links steht Geometrie. Rechts steht Energie und Impuls. Die Gleichung sagt: Materie und Strahlung bestimmen die Raumzeitgeometrie, und die Geometrie bestimmt die Bewegung von Materie und Strahlung.

\begin{merkbox}
Klassische Gravitation ist dynamische Geometrie. Sie ist nicht nur eine Kraft im Raum, sondern eine Dynamik der Raumzeitstruktur selbst.
\end{merkbox}

\chapter{Variationsprinzip und Wirkung}
\section{Das Prinzip der stationären Wirkung}
In der klassischen Physik werden viele Grundgleichungen aus einem Wirkungsfunktional gewonnen. Für Felder schreibt man
\[
S[\phi]=\int \Lag(\phi,\partial\phi,x)\,\dd^4x.
\]

Die physikalischen Gleichungen folgen aus
\[
\delta S=0.
\]

\section{Einstein-Hilbert-Wirkung}
Die gravitative Wirkung lautet
\[
S_G=\frac{c^3}{16\pi G}\int R\sqrt{-g}\,\dd^4x.
\]

Hier bezeichnet $g=\det(g_{\mu\nu})$. Die Variation nach der Metrik liefert die Einstein-Gleichungen, wenn die Materiewirkung hinzugenommen wird:
\[
S = S_G + S_M.
\]

\section{Energie-Impuls-Tensor aus Variation}
Der Energie-Impuls-Tensor ergibt sich aus der Variation der Materiewirkung nach der Metrik:
\[
T_{\mu\nu}=-\frac{2}{\sqrt{-g}}\frac{\delta S_M}{\delta g^{\mu\nu}}.
\]

Das ist wichtig: Materie koppelt nicht beliebig an Gravitation. Sie koppelt über ihre Energie-Impuls-Struktur an die Metrik.

\begin{satz}[Quelle der klassischen Gravitation]
In der Allgemeinen Relativitätstheorie ist nicht Masse allein die Quelle der Gravitation, sondern der gesamte Energie-Impuls-Tensor.
\end{satz}

\begin{proof}
Die Feldgleichung enthält auf der rechten Seite $T_{\mu\nu}$, nicht nur eine skalare Masse. Druck, Energiefluss, Impulsdichte und Spannung tragen ebenfalls zur Quelle bei. Dies folgt direkt aus der kovarianten Kopplung der Materiewirkung an die Metrik.
\end{proof}

\chapter{Quantenprinzip}
\section{Zustandsvektoren}
In der Quantenmechanik wird ein reiner Zustand durch einen Vektor $\ket{\psi}$ in einem Hilbertraum $\Hh$ beschrieben. Physikalisch gleiche Zustände unterscheiden sich höchstens durch einen komplexen Phasenfaktor.

\section{Superposition}
Wenn $\ket{a}$ und $\ket{b}$ mögliche Zustände sind, dann ist auch
\[
\ket{\psi}=\alpha\ket{a}+\beta\ket{b}
\]
mit $\alpha,\beta\in\C$ ein möglicher Zustand, sofern
\[
|\alpha|^2+|\beta|^2=1.
\]

Das Superpositionsprinzip ist kein Randphänomen. Es ist der Kern der Quantenmechanik.

\section{Observable}
Messbare Größen werden durch selbstadjungierte Operatoren beschrieben. Für eine Observable $A$ gilt
\[
A=A^\dagger.
\]

Der Erwartungswert im Zustand $\ket{\psi}$ lautet
\[
\langle A\rangle_\psi = \bra{\psi}A\ket{\psi}.
\]

\section{Messstatistik}
Wenn $A$ Eigenzustände $\ket{a_i}$ mit Eigenwerten $a_i$ besitzt, dann liefert eine Messung mit Wahrscheinlichkeit $|\braket{a_i}{\psi}|^2$ den Wert $a_i$.

Quantenmechanik beschreibt also nicht nur Unkenntnis über klassische Werte. Sie beschreibt Zustände, in denen mehrere klassisch unterscheidbare Möglichkeiten kohärent überlagert sind.

\chapter{Quantenfelder und Energie-Impuls}
\section{Felder werden Operatoren}
In der Quantenfeldtheorie werden Felder zu Operatoren. Ein klassisches Skalarfeld $\phi(x)$ wird zu einem Operatorfeld $\hat\phi(x)$.

Auch der Energie-Impuls-Tensor wird zum Operator:
\[
T_{\mu\nu}(x) \longrightarrow \hat T_{\mu\nu}(x).
\]

Damit besitzt die Quelle der Gravitation quantenmechanische Eigenschaften.

\section{Superposition von Quellen}
Betrachten wir zwei lokalisierte Materiezustände $\ket{L}$ und $\ket{R}$, deren Energieverteilungen unterschiedlich liegen. Eine Quantenüberlagerung lautet
\[
\ket{\psi}=\frac{1}{\sqrt2}(\ket{L}+\ket{R}).
\]

Dann ist die Quelle der Gravitation nicht einfach klassisch entweder links oder rechts. Der Energie-Impuls-Tensor ist ein Operator, und Erwartungswerte können Interferenzstruktur enthalten.

\section{Das Kernproblem}
Klassische Gravitation erwartet auf der rechten Seite der Einstein-Gleichung einen klassischen Tensor $T_{\mu\nu}$. Quantenmaterie liefert aber einen Operator $\hat T_{\mu\nu}$ oder Messstatistiken aus Quantenzuständen.

Eine naheliegende Zwischenlösung wäre
\[
G_{\mu\nu}=\frac{8\pi G}{c^4}\langle\psi|\hat T_{\mu\nu}|\psi\rangle.
\]

Diese Gleichung heißt semiklassische Kopplung. Sie ist als Näherung nützlich, aber nicht als vollständige Endlösung.

\chapter{Warum semiklassische Gravitation nicht genügt}
\section{Semiklassische Kopplung}
Die semiklassische Gleichung lautet vereinfacht
\[
G_{\mu\nu}[g]=\frac{8\pi G}{c^4}\langle \hat T_{\mu\nu}\rangle.
\]

Links steht eine klassische Metrik. Rechts steht der Erwartungswert eines Quantenoperators.

\section{Problem 1: Erwartungswert ist nicht Ereignisstruktur}
Ein Erwartungswert ist nicht dasselbe wie ein einzelnes Messergebnis. Wenn ein Objekt in einer Superposition links und rechts ist, dann kann der Erwartungswert einer gemittelten Quelle entsprechen, obwohl einzelne Messungen lokalisierte Ergebnisse liefern.

Die klassische Metrik wäre dann an eine gemittelte Quelle gekoppelt. Messungen der Materie liefern aber einzelne Ergebnisse. Eine vollständige Theorie muss erklären, wie Metrik, Materiezustand und Messergebnis zusammen konsistent bleiben.

\section{Problem 2: Nichtlinearität}
Die Einstein-Gleichungen sind nichtlinear in der Metrik. Erwartungswerte quantenmechanischer Quellen in eine nichtlineare klassische Gleichung einzusetzen, ist nicht dasselbe wie eine lineare Quantendynamik für das Gesamtsystem.

Nichtlineare Modifikationen der Quantenmechanik sind gefährlich, weil sie in vielen Kontexten zu Problemen mit Superposition, Verschränkung und Signalstruktur führen können.

\section{Problem 3: Rückwirkung}
Materie beeinflusst Geometrie. Geometrie beeinflusst Materie. Wenn Materie quantisiert ist, dann wird die Rückwirkung auf die Geometrie von quantenmechanischen Zuständen abhängen.

Wenn die Geometrie aber strikt klassisch bleibt, entsteht eine hybride Dynamik: ein Teil der Welt ist quantisiert, der andere Teil klassisch, aber beide wechselwirken dynamisch.

Solche hybriden Theorien sind nur unter starken Zusatzbedingungen konsistent. Als allgemeine Endbeschreibung sind sie nicht ausreichend gesichert.

\begin{satz}[Unzulänglichkeit der semiklassischen Endbeschreibung]
Eine Gleichung, die die klassische Metrik ausschließlich an den Erwartungswert des quantisierten Energie-Impuls-Tensors koppelt, kann als Näherung nützlich sein, ersetzt aber keine vollständige konsistente Theorie der wechselwirkenden quantisierten Materie und Gravitation.
\end{satz}

\begin{proof}
Die Gleichung verwendet auf der Materieseite Quantenstatistik, auf der Geometrieseite aber ein einzelnes klassisches Feld. Sie enthält keine gravitativen Quantenfluktuationen, keine Superposition geometrischer Freiheitsgrade und keine vollständige gemeinsame Messstruktur von Materie und Geometrie. Damit kann sie nicht die allgemeinste Dynamik eines Systems beschreiben, in dem die Quelle der Geometrie selbst quantenmechanisch ist.
\end{proof}

\chapter{Schwache Feldnäherung}
\section{Störung um flache Raumzeit}
Für schwache Gravitationsfelder schreibt man
\[
g_{\mu\nu}=\eta_{\mu\nu}+h_{\mu\nu},
\]
wobei $\eta_{\mu\nu}$ die Minkowski-Metrik ist und $h_{\mu\nu}$ eine kleine Störung.

Diese Störung beschreibt Gravitationswellen und schwache gravitative Wechselwirkungen.

\section{Linearisierte Eichsymmetrie}
Unter infinitesimalen Koordinatentransformationen verändert sich $h_{\mu\nu}$ ähnlich wie ein Eichfeld:
\[
h_{\mu\nu}\mapsto h_{\mu\nu}+\partial_\mu\xi_\nu+\partial_\nu\xi_\mu.
\]

Das zeigt: Die schwache Gravitationsstörung besitzt eine Eichredundanz. Nicht jede Komponente von $h_{\mu\nu}$ ist physikalisch unabhängig.

\section{Gravitationswellen}
Im Vakuum erfüllen geeignete Komponenten der Störung eine Wellengleichung:
\[
\Box h_{\mu\nu}=0
\]
unter passenden Eichbedingungen.

Damit besitzt Gravitation im schwachen Feld propagierende Freiheitsgrade.

\begin{satz}[Propagierende gravitative Freiheitsgrade]
Die Allgemeine Relativitätstheorie enthält im schwachen Feld Wellenlösungen. Gravitation besitzt daher dynamische Freiheitsgrade, nicht nur statische Zwangsbedingungen.
\end{satz}

\begin{proof}
Die linearisierte Feldgleichung führt im Vakuum auf eine Wellengleichung für die transvers-traceless Komponenten von $h_{\mu\nu}$. Wellengleichungen beschreiben propagierende Freiheitsgrade. Also besitzt das Gravitationsfeld eigene Dynamik.
\end{proof}

\chapter{Spin-2-Feld und Allgemeine Relativität}
\section{Masselose Felder}
In der relativistischen Quantenfeldtheorie werden freie Teilchenzustände durch Masse und Spin klassifiziert. Ein masseloses Spin-2-Feld besitzt genau die Transformationsstruktur, die zur linearen Gravitationsstörung passt.

\section{Kopplung an Energie-Impuls}
Ein masseloses Spin-2-Feld $h_{\mu\nu}$ koppelt natürlich an den Energie-Impuls-Tensor:
\[
\Lag_{\text{int}} = \kappa h_{\mu\nu}T^{\mu\nu}.
\]

Hier ist $\kappa$ eine Kopplungskonstante, die mit Newtons Gravitationskonstante zusammenhängt.

\section{Selbstkopplung}
Das Gravitationsfeld trägt selbst Energie und Impuls. Daher muss es auch an seine eigene Energie koppeln. Diese Forderung führt von der linearen Theorie schrittweise zur nichtlinearen Struktur der Allgemeinen Relativitätstheorie.

\begin{satz}[Niederenergetische Eindeutigkeit]
Eine konsistente relativistische Theorie eines masselosen Spin-2-Feldes, das universell an den Energie-Impuls-Tensor koppelt, reproduziert bei niedrigen Energien die Struktur der Allgemeinen Relativitätstheorie.
\end{satz}

\begin{proof}
Die Eichsymmetrie des masselosen Spin-2-Feldes erzwingt die Erhaltung der gekoppelten Quelle. Da das Feld selbst Energie trägt, muss die Quelle die Energie des Feldes einschließen. Iteriert man diese Selbstkopplung konsistent, entsteht die nichtlineare kovariante Struktur, die durch die Einstein-Gleichungen beschrieben wird. Dieser Standardweg zeigt die niederenergetische Besonderheit der Einstein-Gravitation.
\end{proof}

\chapter{Warum das Gravitationsfeld quantisiert werden muss}
\section{Allgemeines Argument}
Wenn ein Feld dynamische Freiheitsgrade besitzt und mit Quantenmaterie wechselwirkt, dann kann es Information über Quantenzustände tragen. Ein klassisches Feld, das mit einer Quantenquelle wechselwirkt, würde entweder Superpositionen zerstören, nur Erwartungswerte sehen oder eine hybride Dynamik erzeugen.

Keine dieser drei Möglichkeiten ist als allgemeine Endbeschreibung befriedigend.

\section{Superpositionsargument}
Betrachten wir eine Masse in Superposition:
\[
\ket{\psi}=\frac{1}{\sqrt2}(\ket{L}+\ket{R}).
\]

Wenn die Gravitation klassisch wäre, müsste es ein einzelnes klassisches Gravitationsfeld geben. Dieses Feld müsste entweder zur linken, zur rechten oder zur gemittelten Verteilung gehören.

Alle drei Möglichkeiten sind problematisch:

\begin{enumerate}[label=\arabic*.]
\item Links oder rechts würde die Symmetrie der Superposition ohne Messung brechen.
\item Das Mittel würde eine Quelle beschreiben, die keiner einzelnen Messrealisation entspricht.
\item Eine stochastische Auswahl müsste zusätzliche Kollapsdynamik enthalten.
\end{enumerate}

Eine quantisierte Gravitation erlaubt dagegen einen verschränkten Zustand der Form
\[
\ket{\Psi}=\frac{1}{\sqrt2}(\ket{L}\ket{g_L}+\ket{R}\ket{g_R}).
\]

Hier trägt das Gravitationsfeld selbst quantenmechanische Alternativen.

\section{Verschränkungsargument}
Ein klassischer Vermittler kann unter Standardannahmen keine Verschränkung zwischen zwei Quantensystemen erzeugen. Wenn zwei Massen nur über Gravitation wechselwirken und danach verschränkt sind, dann muss der Vermittler nichtklassische Freiheitsgrade besitzen.

\begin{satz}[Nichtklassischer Vermittler]
Wenn zwei Quantensysteme, die nicht direkt wechselwirken, durch eine vermittelnde Wechselwirkung verschränkt werden, dann muss der Vermittler selbst nichtklassische Freiheitsgrade besitzen, sofern die Dynamik lokal und ohne versteckte direkte Kopplung beschrieben wird.
\end{satz}

\begin{proof}
Rein klassische Kommunikation zwischen Quantensystemen kann Korrelationen erzeugen, aber keine Verschränkung aus zunächst separablen Zuständen unter lokalen Operationen und klassischer Kommunikation. Entsteht Verschränkung ausschließlich über den Vermittler, dann kann dieser Vermittler nicht nur ein klassischer Informationskanal sein. Auf Gravitation angewandt bedeutet dies: Kann Gravitation Verschränkung zwischen Massen vermitteln, besitzt sie nichtklassische Freiheitsgrade.
\end{proof}

\chapter{Der niederenergetische Quantengravitationssatz}
\section{Quantisierung der linearen Störung}
Die lineare Gravitationsstörung $h_{\mu\nu}$ kann wie andere freie Felder in Moden zerlegt werden. Für jede physikalische Mode entstehen Erzeugungs- und Vernichtungsoperatoren.

Formal:
\[
\hat h_{\mu\nu}(x)=\sum_\lambda\int \frac{\dd^3k}{(2\pi)^3}\,\mathcal{N}_k
\left(\epsilon^{(\lambda)}_{\mu\nu}(k)\hat a_{k,\lambda}e^{-ikx}+\epsilon^{(\lambda)*}_{\mu\nu}(k)\hat a^\dagger_{k,\lambda}e^{ikx}\right).
\]

Die Quanten dieses Feldes heißen Gravitonen. In der niederenergetischen Theorie sind sie masselose Spin-2-Anregungen.

\section{Effektive Feldtheorie}
Da Newtons Konstante eine dimensionsbehaftete Kopplung ist, ist die Quantentheorie der Gravitation perturbativ nicht renormierbar im strengen Sinn einer einfachen fundamentalen Theorie. Sie ist aber als effektive Feldtheorie bei niedrigen Energien wohldefiniert.

Die Wirkung enthält dann
\[
S=\int \sqrt{-g}\left(\frac{c^3}{16\pi G}R + c_1 R^2 + c_2 R_{\mu\nu}R^{\mu\nu}+\cdots\right)\dd^4x.
\]

Höhere Krümmungsterme werden bei niedrigen Energien unterdrückt.

\section{Der Satz}
\begin{satz}[Niederenergetische Quantengravitation]
Unter den Annahmen der Lorentz-Invarianz, lokalen Energie-Impuls-Erhaltung, universellen Kopplung an Energie und Impuls, quantisierter Materie und dynamischer schwacher Gravitationsfelder ist die niederenergetische Theorie der Gravitation eine Quantentheorie masseloser Spin-2-Freiheitsgrade, deren klassische Grenze die Allgemeine Relativitätstheorie ist.
\end{satz}

\begin{proof}
Die Allgemeine Relativitätstheorie besitzt in der schwachen Feldnäherung propagierende Freiheitsgrade $h_{\mu\nu}$. Diese transformieren unter linearisierten Koordinatentransformationen wie ein masseloses Spin-2-Feld mit Eichredundanz. Die Kopplung an Materie erfolgt universell über den Energie-Impuls-Tensor. Da Materie quantisiert ist, ist $T_{\mu\nu}$ operatorwertig und kann Superpositionen und Verschränkungen tragen. Eine rein klassische Behandlung von $h_{\mu\nu}$ kann diese Struktur nicht allgemein ohne Verlust von Superposition, ohne bloße Erwartungswertkopplung oder ohne hybride Zusatzdynamik beschreiben. Daher muss $h_{\mu\nu}$ selbst als Quantenfeld behandelt werden. Die nichtlineare Selbstkopplung reproduziert in der klassischen Grenze die Einstein-Struktur. Da die Kopplung dimensionsbehaftet ist, ist die niederenergetische Form eine effektive Feldtheorie. Damit folgt die niederenergetische Quantengravitation als notwendige Konsequenz der genannten Prinzipien.
\end{proof}

\chapter{Planck-Skala}
\section{Dimensionsanalyse}
Aus $G$, $\hbar$ und $c$ lassen sich charakteristische Planck-Größen bilden.

Planck-Länge:
\[
\ell_P=\sqrt{\frac{\hbar G}{c^3}}.
\]

Planck-Zeit:
\[
t_P=\sqrt{\frac{\hbar G}{c^5}}.
\]

Planck-Masse:
\[
m_P=\sqrt{\frac{\hbar c}{G}}.
\]

\section{Bedeutung}
Die Planck-Skala markiert den Bereich, in dem Quantenwirkung und gravitative Wirkung gleichzeitig stark werden.

Das bedeutet nicht automatisch, dass Raum aus kleinen Würfeln besteht. Es bedeutet: Die bekannten klassischen und perturbativen Beschreibungen verlieren dort ihre einfache Gültigkeit.

\section{Grenze des Skripts}
Dieses Skript beweist die Notwendigkeit quantisierter gravitativer Freiheitsgrade und beschreibt die niederenergetische Quantengravitation. Die vollständige Planck-Skalen-Theorie bleibt eine weitergehende Aufgabe.

\chapter{Schwarze Löcher und Thermodynamik}
\section{Horizonte}
Ein schwarzes Loch besitzt einen Ereignishorizont. Für ein nichtrotierendes, ungeladenes schwarzes Loch lautet der Schwarzschild-Radius
\[
r_s=\frac{2GM}{c^2}.
\]

\section{Bekenstein-Hawking-Entropie}
Die Entropie eines schwarzen Lochs ist proportional zur Horizontfläche:
\[
S_{BH}=\frac{k_B c^3 A}{4G\hbar}.
\]

Diese Formel enthält zugleich $G$, $c$ und $\hbar$. Sie ist daher ein starkes Indiz dafür, dass Gravitation, Quantenphysik und Thermodynamik gemeinsam beschrieben werden müssen.

\section{Hawking-Temperatur}
Die Hawking-Temperatur lautet
\[
T_H=\frac{\hbar c^3}{8\pi G M k_B}.
\]

Auch hier treten Gravitationskonstante und Plancksches Wirkungsquantum gemeinsam auf.

\begin{satz}[Thermodynamischer Hinweis]
Die Thermodynamik schwarzer Löcher zeigt, dass eine vollständige Theorie der Gravitation Quanteneigenschaften enthalten muss.
\end{satz}

\begin{proof}
Entropie und Temperatur schwarzer Löcher enthalten $\hbar$. Eine rein klassische Theorie enthält dieses Wirkungsquantum nicht als dynamische Grundlage. Da Horizonte gravitative Objekte sind, aber thermische Quanteneigenschaften zeigen, ist eine Verbindung von Gravitation, Quantenphysik und Thermodynamik erforderlich.
\end{proof}

\chapter{Konsistenz der Beweiskette}
\section{Die Kette}
Die Beweiskette lässt sich knapp schreiben:

\[
\text{Quantenmaterie}\Rightarrow \hat T_{\mu\nu}
\]

\[
\hat T_{\mu\nu}\text{ koppelt an }g_{\mu\nu}
\]

\[
g_{\mu\nu}=\eta_{\mu\nu}+h_{\mu\nu}
\]

\[
h_{\mu\nu}\text{ besitzt propagierende Freiheitsgrade}
\]

\[
\text{propagierende wechselwirkende Freiheitsgrade mit Quantenquellen müssen quantisiert werden}
\]

\[
\Rightarrow \hat h_{\mu\nu}
\]

\[
\Rightarrow \text{niederenergetische Quantengravitation.}
\]

\section{Was genau bewiesen ist}
Bewiesen ist:

\begin{enumerate}[label=\arabic*.]
\item Gravitation besitzt dynamische Feldfreiheitsgrade.
\item Diese Freiheitsgrade koppeln an Energie und Impuls.
\item Energie und Impuls quantisierter Materie sind nicht allgemein klassische Größen.
\item Eine rein klassische Geometrie kann Superposition, Rückwirkung und Verschränkung nicht allgemein vollständig tragen.
\item Die schwache Gravitationsstörung ist ein masseloses Spin-2-Feld.
\item Dieses Feld muss bei konsistenter Kopplung an Quantenmaterie quantisiert werden.
\item Die niederenergetische Form ist eine effektive Quantentheorie der Gravitation.
\end{enumerate}

\section{Was nicht bewiesen ist}
Nicht bewiesen ist:

\begin{enumerate}[label=\arabic*.]
\item Welche Theorie bei beliebig hohen Energien endgültig gilt.
\item Ob Raumzeit fundamental diskret, emergent oder kontinuierlich ist.
\item Ob die vollständige Theorie durch Schleifen, Strings, Holographie, kausale Mengen oder eine andere Struktur beschrieben wird.
\item Ob alle Singularitäten ohne Zusatzannahmen aufgelöst werden.
\end{enumerate}

Diese Begrenzung schwächt den Beweis nicht. Sie macht ihn präzise.

\chapter{Ausblick auf bekannte Forschungswege}
\section{Effektive Quantengravitation}
Die effektiv-feldtheoretische Sicht behandelt Allgemeine Relativität als niederenergetische Theorie. Sie erlaubt berechenbare Quantenkorrekturen bei Energien weit unterhalb der Planck-Skala.

\section{Schleifenquantengravitation}
Die Schleifenquantengravitation versucht, geometrische Größen wie Fläche und Volumen direkt zu quantisieren.

\section{Stringtheorie}
Die Stringtheorie ersetzt punktförmige Teilchen durch eindimensionale Objekte. Ein masseloser Spin-2-Zustand erscheint dort natürlich im Spektrum.

\section{Kausale Mengen}
Kausale-Mengen-Ansätze beschreiben Raumzeit fundamental als diskrete kausale Ordnung.

\section{Holographie}
Holographische Ansätze legen nahe, dass gravitative Information in bestimmten Fällen auf Randtheorien abgebildet werden kann.

Diese Wege sind verschieden, aber sie teilen den Grundbefund: Klassische Raumzeit allein ist nicht die letzte Beschreibung.


\chapter{Hamiltonsche Sicht auf die Raumzeit}
\section{Warum eine Hamiltonsche Form nützlich ist}
Die Lagrange-Formulierung beschreibt die Dynamik aus der Wirkung. Für die Quantisierung ist zusätzlich die Hamiltonsche Sicht nützlich, weil Quantentheorie häufig von Zuständen und ihren zeitlichen Entwicklungen ausgeht.

In der Mechanik beginnt man mit Koordinaten $q^i$ und Geschwindigkeiten $\dot q^i$. Die kanonischen Impulse sind
\[
p_i=\frac{\partial L}{\partial \dot q^i}.
\]
Die Hamilton-Funktion lautet
\[
H(q,p)=p_i\dot q^i-L(q,\dot q).
\]

Bei Feldern wird aus dieser Konstruktion eine Feldtheorie auf Raumhyperflächen. In der Gravitation ist das besonders anspruchsvoll, weil die Raumzeit selbst dynamisch ist und daher die Zerlegung in Raum und Zeit nicht nur eine technische Wahl ist.

\section{ADM-Zerlegung}
In der ADM-Formulierung wird die Raumzeit in räumliche Hyperflächen zerlegt. Die Metrik wird durch drei Bestandteile beschrieben:

\begin{enumerate}[label=\arabic*.]
\item die räumliche Metrik $\gamma_{ij}$,
\item die Lapse-Funktion $N$,
\item den Shift-Vektor $N^i$.
\end{enumerate}

Anschaulich beschreibt $\gamma_{ij}$ die Geometrie einer Raumfläche, $N$ den Abstand zwischen benachbarten Flächen und $N^i$ die seitliche Verschiebung der Koordinaten.

\section{Zwangsbedingungen}
Die Allgemeine Relativitätstheorie besitzt nicht nur Bewegungsgleichungen, sondern auch Zwangsbedingungen. Sie entstehen aus der Diffeomorphismeninvarianz der Theorie. Physikalische Zustände dürfen nicht von bloßer Koordinatenwahl abhängen.

Die Hamiltonsche Nebenbedingung und die Impulsnebenbedingungen drücken genau das aus. In symbolischer Form:
\[
\mathcal{H}=0,\qquad \mathcal{H}_i=0.
\]

Diese Gleichungen sind kein mathematischer Schmuck. Sie sagen: Nicht jede formale Konfiguration ist ein physikalisch zulässiger Zustand.

\section{Bedeutung für Quantisierung}
Bei der Quantisierung werden klassische Größen zu Operatoren. Aus Nebenbedingungen werden Bedingungen an zulässige Quantenzustände:
\[
\hat{\mathcal{H}}\ket{\Psi}=0,\qquad \hat{\mathcal{H}}_i\ket{\Psi}=0.
\]

Damit erscheint die Quantengravitation nicht nur als Quantisierung eines weiteren Feldes, sondern auch als Quantisierung einer Theorie mit geometrischen Zwangsbedingungen.

\chapter{Kanonische Quantisierung}
\section{Vom klassischen Poisson-Klammerausdruck zum Kommutator}
In der kanonischen Quantisierung ersetzt man klassische Poisson-Klammern durch Kommutatoren:
\[
\{q^i,p_j\}=\delta^i_j
\quad\longrightarrow\quad
[\hat q^i,\hat p_j]=i\hbar\delta^i_j.
\]

Für Felder entsteht entsprechend
\[
[\hat\phi(t,\mathbf{x}),\hat\pi(t,\mathbf{y})]=i\hbar\delta^{(3)}(\mathbf{x}-\mathbf{y}).
\]

\section{Anwendung auf die Metrik}
In einer kanonischen Gravitationstheorie wird die räumliche Metrik $\gamma_{ij}$ zur Konfigurationsgröße. Ihr kanonischer Impuls hängt mit der extrinsischen Krümmung zusammen. Formal führt die Quantisierung zu Operatoren
\[
\hat\gamma_{ij},\qquad \hat\pi^{ij}.
\]

Die zentrale Schwierigkeit besteht darin, dass diese Operatoren nicht auf einem festen Hintergrundraum wirken, sondern die Geometrie selbst betreffen.

\section{Wheeler-DeWitt-Gleichung}
Eine symbolische Form der kanonischen Quantengravitation ist die Wheeler-DeWitt-Gleichung:
\[
\hat{\mathcal{H}}\Psi[\gamma]=0.
\]

Hier ist $\Psi[\gamma]$ ein Funktional der räumlichen Geometrie. Diese Gleichung zeigt in maximal knapper Form, wie radikal der Schritt ist: Nicht ein Feld auf Raumzeit wird quantisiert, sondern geometrische Struktur selbst wird Teil des Quantenzustands.

\section{Warum das den Hauptsatz stützt}
Die kanonische Sicht liefert keinen einfachen fertigen Rechenapparat für alle Situationen. Sie bestätigt aber den zentralen Zwang: Wenn die Metrik dynamisch ist und ihre Freiheitsgrade echte physikalische Freiheitsgrade sind, dann führt eine Quantisierung nicht um die Geometrie herum. Sie betrifft die Geometrie selbst.

\chapter{Pfadintegral und Summation über Geometrien}
\section{Pfadintegral in der Quantenmechanik}
In der Pfadintegralformulierung erhält man Übergangsamplituden durch Summation über Wege:
\[
\langle x_f,t_f|x_i,t_i\rangle=\int \mathcal{D}x\,e^{\frac{i}{\hbar}S[x]}.
\]

Klassische Bahnen erscheinen als stationäre Beiträge im Grenzfall $\hbar\to 0$.

\section{Feldtheoretisches Pfadintegral}
Für ein Feld schreibt man formal
\[
Z=\int \mathcal{D}\phi\,e^{\frac{i}{\hbar}S[\phi]}.
\]

Das Pfadintegral ist nicht immer mathematisch streng definiert, aber es ist ein sehr wirksames Werkzeug der Quantenfeldtheorie.

\section{Gravitativer Fall}
Für Gravitation liegt formal nahe:
\[
Z=\int \mathcal{D}g\,e^{\frac{i}{\hbar}S[g]}.
\]

Das bedeutet: Es wird nicht nur über Materiefelder auf fester Raumzeit summiert, sondern über Metriken selbst. Damit wird Raumzeitgeometrie in die Quantensumme aufgenommen.

\section{Klassische Grenze}
Im Grenzfall großer Wirkung gegenüber $\hbar$ dominiert die stationäre Phase. Dann gilt näherungsweise
\[
\delta S[g]=0,
\]
und daraus folgen die klassischen Einstein-Gleichungen.

Damit zeigt die Pfadintegralformulierung denselben Zusammenhang von der anderen Seite: Klassische Raumzeit erscheint als Grenzfall einer quantenhaften Summation über geometrische Möglichkeiten.

\chapter{Klassische Grenze und Korrespondenz}
\section{Warum die klassische Theorie erhalten bleiben muss}
Eine Quantengravitation darf die Allgemeine Relativitätstheorie nicht einfach ersetzen, als wäre diese falsch. Sie muss sie in ihrem Gültigkeitsbereich reproduzieren.

Das nennt man Korrespondenzprinzip. Eine neue Theorie muss die bewährte alte Theorie als Grenzfall enthalten.

\section{Der Grenzübergang}
In geeigneten Zuständen muss gelten:
\[
\langle \hat g_{\mu\nu}\rangle \approx g_{\mu\nu}^{\text{klassisch}}.
\]

Für schwache Felder muss außerdem die Quantentheorie des Spin-2-Feldes in den klassischen Grenzfall von Gravitationswellen übergehen.

\section{Drei Grenzfälle}
Eine tragfähige Theorie muss mindestens drei Grenzfälle besitzen:

\begin{enumerate}[label=\arabic*.]
\item Für $G\to 0$ muss gewöhnliche Quantenfeldtheorie ohne dynamische Gravitation erscheinen.
\item Für $\hbar\to 0$ muss klassische Allgemeine Relativität erscheinen.
\item Für schwache Felder und niedrige Energien muss die effektive Quantentheorie des masselosen Spin-2-Feldes erscheinen.
\end{enumerate}

Diese Grenzfälle sind Prüfsteine jeder Kandidatentheorie.

\chapter{Beobachtungsbezug}
\section{Warum ein mathematischer Zwang noch keine Messung ersetzt}
Der Beweis der Notwendigkeit quantisierter Gravitation zeigt, dass eine rein klassische Endbeschreibung unzureichend ist. Eine Messung konkreter Gravitationsquanten oder gravitationsvermittelter Verschränkung ist ein zusätzlicher Schritt.

Mathematik zeigt hier die Richtung. Experimente entscheiden über konkrete Modelle und Parameter.

\section{Gravitationswellen}
Die direkte Beobachtung von Gravitationswellen bestätigt, dass gravitative Freiheitsgrade propagieren. Sie bestätigt noch nicht einzelne Gravitonen, aber sie bestätigt den dynamischen Feldcharakter der Gravitation.

\section{Kosmologie}
Frühuniverselle Fluktuationen, Hintergrundstrahlung und Strukturbildung können Hinweise auf quantisierte Anfangsbedingungen enthalten. Die Interpretation ist anspruchsvoll, weil kosmologische Daten Modellannahmen benötigen.

\section{Laborexperimente mit Massen in Superposition}
Vorgeschlagene Experimente prüfen, ob Gravitation zwischen kleinen Massen Verschränkung erzeugen kann. Ein positives Ergebnis wäre ein starkes Indiz für nichtklassische gravitative Freiheitsgrade.

\section{Schwarze Löcher}
Schwarze Löcher verbinden Gravitation, Quantenphysik und Thermodynamik besonders eng. Hawking-Strahlung und Horizontentropie sind deshalb zentrale Hinweise auf die Unvollständigkeit rein klassischer Gravitation.

\chapter{Didaktische Zusammenfassung der Beweise}
\section{Erster Beweisweg: Quelle}
\begin{enumerate}[label=\arabic*.]
\item Gravitation koppelt an Energie und Impuls.
\item Energie und Impuls quantisierter Materie sind operatorwertig.
\item Eine klassische Metrik kann einen Operator nicht vollständig als klassische Quelle aufnehmen.
\item Daher ist eine rein klassische Endbeschreibung unzureichend.
\end{enumerate}

\section{Zweiter Beweisweg: Feld}
\begin{enumerate}[label=\arabic*.]
\item Die Metrik besitzt im schwachen Feld propagierende Störungen.
\item Diese Störungen haben die Struktur eines masselosen Spin-2-Feldes.
\item Wechselwirkende Felder mit Quantenquellen müssen quantisiert werden.
\item Daher muss das gravitative Feld quantisierte Freiheitsgrade besitzen.
\end{enumerate}

\section{Dritter Beweisweg: Verschränkung}
\begin{enumerate}[label=\arabic*.]
\item Ein klassischer Vermittler erzeugt unter Standardannahmen keine Verschränkung zwischen getrennten Quantensystemen.
\item Gravitation soll zwischen Massen vermitteln.
\item Wenn Gravitation Verschränkung vermitteln kann, muss sie nichtklassische Freiheitsgrade besitzen.
\item Damit ist eine rein klassische Gravitation ausgeschlossen.
\end{enumerate}

\section{Vierter Beweisweg: Thermodynamik}
\begin{enumerate}[label=\arabic*.]
\item Schwarze Löcher sind gravitative Objekte.
\item Ihre Entropie und Temperatur enthalten $\hbar$.
\item Also ist ihre vollständige Beschreibung quantenphysikalisch.
\item Damit reicht klassische Gravitation nicht aus.
\end{enumerate}

\section{Zusammenführung}
Alle vier Wege führen zur gleichen Aussage: Die Gravitation kann in ihrem vollständigen Wechselwirkungszusammenhang nicht rein klassisch bleiben.

\chapter{Übungen}
\section*{Übung 1: Metrik und Kausalstruktur}
Erklären Sie, warum eine dynamische Metrik zugleich eine dynamische Kausalstruktur bedeutet.

\section*{Übung 2: Energie-Impuls-Tensor}
Warum reicht Masse allein nicht als Quelle der Gravitation?

\section*{Übung 3: Superposition}
Betrachten Sie den Zustand
\[
\ket{\psi}=\frac{1}{\sqrt2}(\ket{L}+\ket{R}).
\]
Warum ist eine gemittelte klassische Quelle keine vollständige Beschreibung aller Messsituationen?

\section*{Übung 4: Linearisierte Gravitation}
Zeigen Sie aus
\[
g_{\mu\nu}=\eta_{\mu\nu}+h_{\mu\nu}
\]
warum $h_{\mu\nu}$ als Feldstörung verstanden werden kann.

\section*{Übung 5: Spin-2}
Warum ist die universelle Kopplung an $T_{\mu\nu}$ ein Hinweis auf Spin 2 und nicht auf ein skalares Feld?

\section*{Übung 6: Planck-Länge}
Leiten Sie durch Dimensionsanalyse
\[
\ell_P=\sqrt{\frac{\hbar G}{c^3}}
\]
her.

\section*{Übung 7: Beweisgrenze}
Formulieren Sie in eigenen Worten, warum dieses Skript die Notwendigkeit quantisierter Gravitation beweist, aber nicht automatisch eine vollständige Planck-Skalen-Theorie auswählt.

\chapter{Musterlösungen}
\section*{Lösung 1}
Die Metrik entscheidet, ob ein Vektor zeitartig, lichtartig oder raumartig ist. Damit entscheidet sie, welche Ereignisse kausal verbunden sein können. Wenn die Metrik dynamisch ist, ist auch diese Kausalstruktur dynamisch.

\section*{Lösung 2}
In der Allgemeinen Relativitätstheorie steht auf der rechten Seite der Feldgleichung der Energie-Impuls-Tensor. Dieser enthält Energie, Impuls, Druck, Spannung und Flüsse. Masse allein ist nur ein Spezialfall.

\section*{Lösung 3}
Eine gemittelte Quelle beschreibt den Erwartungswert. Messungen liefern aber einzelne Ergebnisse. Eine vollständige Theorie muss erklären, wie Geometrie, Messstatistik und Rückwirkung zusammenpassen.

\section*{Lösung 4}
Die Zerlegung trennt die flache Hintergrundmetrik von einer kleinen Störung. Diese Störung kann Wellengleichungen erfüllen und Energie transportieren. Daher besitzt sie Feldcharakter.

\section*{Lösung 5}
Der Energie-Impuls-Tensor ist ein symmetrischer Tensor zweiter Stufe. Ein Feld, das universell und linear daran koppelt, muss dieselbe Tensorstruktur tragen. Das führt zum masselosen Spin-2-Feld.

\section*{Lösung 6}
Gesucht ist eine Länge aus $\hbar$, $G$ und $c$. Mit Dimensionen
\[
[\hbar]=ML^2T^{-1},\quad [G]=M^{-1}L^3T^{-2},\quad [c]=LT^{-1}
\]
ergibt
\[
\frac{\hbar G}{c^3}\sim L^2.
\]
Die Wurzel ist also eine Länge.

\section*{Lösung 7}
Das Skript zeigt, dass gravitative Freiheitsgrade quantisiert werden müssen, wenn sie konsistent mit Quantenmaterie wechselwirken. Es zeigt aber nicht, welche mathematische Struktur bei Planck-Energien endgültig gilt.

\chapter{Schluss}
Die bekannte Physik führt auf einen klaren Zwang: Materie ist quantisiert; Energie und Impuls sind Quellen der Gravitation; Gravitation ist Dynamik der Raumzeit; die schwache Raumzeitstörung besitzt eigene propagierende Freiheitsgrade; diese Freiheitsgrade wechselwirken mit quantisierten Quellen. Daher können sie nicht allgemein und endgültig rein klassisch bleiben.

Damit ist die Notwendigkeit der Quantengravitation bewiesen.

Die vollständige Theorie der Planck-Skala bleibt eine offene Forschungsaufgabe. Aber der Übergang von klassischer Gravitation zu quantisierten gravitativen Freiheitsgraden ist keine bloße Spekulation. Er folgt aus dem Zusammenspiel von Quantenprinzip, Energie-Impuls-Kopplung, Raumzeitdynamik und relativistischer Feldtheorie.

\begin{center}
\textbf{q.e.d.}\\[2mm]
Ingolf Lohmann
\end{center}

\end{document}
